\section*{Bloque D. Ampliación: Teoría y Aplicaciones Prácticas (70\% / 30\%)}

\noindent\textit{Preguntas 91 a 180. Ampliación del banco oficial con 63 preguntas teóricas y conceptuales (91 a 153) y 27 aplicaciones cuantitativas prácticas (154 a 180).}

\subsection*{Parte I. Preguntas teóricas y conceptuales (91 a 153)}

\q{En economía del bienestar, a diferencia de una externalidad tecnológica (real), una externalidad pecuniaria que opera únicamente a través de variaciones en los precios de mercado:}
\begin{opts}
    \item No genera por sí misma una falla de eficiencia en el sentido de Pareto si los mercados son competitivos.
    \item Exige necesariamente la imposición de un impuesto pigouviano equivalente al cambio porcentual en el precio.
    \item Violenta directamente las funciones de producción físicas al modificar la productividad marginal de los insumos.
    \item Impide la existencia de cualquier equilibrio walrasiano competitivo incluso bajo información completa.
\end{opts}

\q{Cuando una industria competitiva genera contaminación atmosférica ($CMS = CMP + CME$), el área del triángulo de Pérdida de Eficiencia Social ($PIE$) en el equilibrio de mercado no regulado mide:}
\begin{opts}
    \item La recaudación tributaria total que obtendría el Estado si aplicara un impuesto pigouviano óptimo.
    \item El exceso acumulado de $CMS$ sobre $BMS$ en el tramo entre $Q_{soc}^*$ y $Q_M$.
    \item El beneficio privado total que pierden los productores cuando se reduce la contaminación al nivel cero.
    \item La diferencia entre el excedente del consumidor en monopolio y el excedente bajo competencia perfecta.
\end{opts}

\q{En un mercado competitivo donde el consumo de un bien genera una externalidad negativa hacia terceros (por ejemplo, congestión o tabaquismo pasivo), se cumple en el equilibrio libre que:}
\begin{opts}
    \item El beneficio marginal social es superior al beneficio marginal privado y el bien resulta subconsumido.
    \item El costo marginal social coincide con el costo marginal privado, por lo que no existe ineficiencia asignativa.
    \item El beneficio marginal social es menor que el privado ($BMS < BMP$) y existe sobreconsumo ($Q_M > Q_{soc}^*$).
    \item El precio de equilibrio de mercado se sitúa por encima del óptimo social debido a la escasez artificial.
\end{opts}

\q{Cuando una actividad productiva genera externalidades positivas (como la conservación forestal en cabeceras de cuenca que mejora la calidad hídrica aguas abajo), el equilibrio competitivo sin intervención implica:}
\begin{opts}
    \item Que el precio de mercado excede la disposición marginal a pagar de todos los consumidores beneficiados.
    \item Un costo marginal social superior al privado ($CMS > CMP$) y una sobreproducción de servicios ambientales.
    \item La eliminación automática de la falla de mercado mediante la libre entrada de empresas en el largo plazo.
    \item Un costo marginal social inferior al privado ($CMS < CMP$) y una subproducción respecto al óptimo social.
\end{opts}

\q{Si una empresa competitiva genera contaminación con un costo marginal externo $CME(Q)$ estrictamente creciente en la producción $Q$, un impuesto pigouviano fijado igual a $CME(Q_M)$ (evaluado en el equilibrio de mercado inicial) resultará:}
\begin{opts}
    \item Excesivo ($t > t^*$), reduciendo la producción por debajo del óptimo social.
    \item Exactamente óptimo porque el daño marginal observado antes de intervenir refleja la externalidad relevante.
    \item Insuficiente ($t < t^*$), dejando un remanente de sobreproducción contaminante superior al óptimo social.
    \item Equivalente a una asignación de permisos negociables gratuitos sin impacto sobre la cantidad producida.
\end{opts}

\q{En economía ambiental y finanzas públicas, la 'hipótesis del doble dividendo' de la tributación ambiental (impuestos verdes) sostiene que gravar emisiones contaminantes puede:}
\begin{opts}
    \item Duplicar la tasa de crecimiento del producto potencial reduciendo a cero la contaminación industrial.
    \item Mejorar el ambiente y reducir distorsiones reciclando ingresos en otros impuestos.
    \item Eliminar la necesidad de medir el daño marginal externo al fijar una tasa uniforme sobre el consumo.
    \item Sustituir integralmente los derechos de propiedad privada por regulación directa de comando y control.
\end{opts}

\q{Al comparar un impuesto pigouviano por unidad de emisión frente a un subsidio por unidad de emisión reducida (abatida) en una industria competitiva con libre entrada a largo plazo:}
\begin{opts}
    \item Ambos generan exactamente el mismo número de empresas y nivel total de contaminación tanto a corto como a largo plazo.
    \item El subsidio al abatimiento es siempre superior en el largo plazo porque reduce los costos medios y atrae más empresas limpias.
    \item En el largo plazo el subsidio induce entrada de firmas y eleva la contaminación total.
    \item El impuesto pigouviano impide que las empresas adopten tecnologías de abatimiento en el corto plazo por falta de liquidez.
\end{opts}

\q{Si un monopolista produce un bien cuya fabricación genera una externalidad negativa de contaminación, el impuesto pigouviano estándar $t = CME(Q_{soc}^*)$:}
\begin{opts}
    \item No tiene ningún efecto sobre la cantidad producida porque el monopolista absorbe íntegramente el impuesto.
    \item Siempre restaura el óptimo de Pareto (First-Best) al eliminar simultáneamente el poder de mercado y la contaminación.
    \item Debe elevarse por encima del $CME(Q_{soc}^*)$ para compensar las rentas extraordinarias del monopolista.
    \item Puede reducir el bienestar social porque el monopolio ya restringe la producción y el impuesto la contrae aún más.
\end{opts}

\q{Aun cuando los costos de transacción sean nulos y los derechos de propiedad estén perfectamente definidos, la asignación final eficiente del Teorema de Coase dejará de ser independiente de quién reciba el derecho inicial si:}
\begin{opts}
    \item Existen efectos riqueza que alteran la disposición marginal a pagar ambiental.
    \item Las partes tienen distinta capacidad de negociación para repartirse el excedente generado por el acuerdo.
    \item El derecho de propiedad se asigna inicialmente al contaminador en lugar de asignarse a la parte afectada.
    \item La externalidad involucra únicamente a dos agentes económicos con información completa sobre sus funciones.
\end{opts}

\q{En un conflicto ambiental donde una empresa negocia con cientos de propietarios afectados, el 'problema de retención' (hold-out) impide alcanzar la solución de Coase porque:}
\begin{opts}
    \item Los propietarios carecen legalmente de derechos de propiedad transmisibles sobre el recurso ambiental afectado.
    \item Cada afectado tiene incentivos estratégicos a retrasar el acuerdo para exigir una mayor porción del excedente.
    \item La empresa contaminadora siempre desconoce su propia función de costos marginales de abatimiento.
    \item Los costos marginales de abatimiento son nulos cuando existen múltiples partes involucradas en la negociación.
\end{opts}

\q{La ventaja de eficiencia en costos (costo-efectividad) de un impuesto a las emisiones o de permisos negociables frente a un estándar uniforme de reducción (comando y control) radica en que los instrumentos de mercado:}
\begin{opts}
    \item Obligan a todas las empresas a reducir exactamente el mismo porcentaje de sus emisiones históricas.
    \item Eliminan por completo las emisiones contaminantes en el primer período sin requerir monitoreo ambiental.
    \item Igualan los costos marginales de abatimiento entre todas las firmas, logrando la meta al mínimo costo total.
    \item Exoneran de cualquier esfuerzo de abatimiento a las empresas más grandes que operan con economías de escala.
\end{opts}

\q{En un sistema competitivo de permisos de emisión negociables (Cap-and-Trade) sin costos de transacción ni distorsiones fiscales, comparar la subasta inicial de permisos frente a su asignación gratuita (grandfathering) muestra que:}
\begin{opts}
    \item La subasta genera mayores emisiones totales porque las empresas compran permisos adicionales para contaminar más.
    \item Solo la subasta logra igualar los costos marginales de abatimiento entre las empresas participantes.
    \item La asignación gratuita impide el comercio secundario de permisos y eleva el precio de equilibrio del permiso.
    \item Ambos métodos alcanzan la misma asignación eficiente de abatimiento, pero difieren en la distribución de rentas.
\end{opts}

\q{Según el teorema de Weitzman (1974) sobre regulación ambiental bajo incertidumbre en los costos marginales de abatimiento, un sistema de cuotas o permisos negociables (cantidades) es preferible a un impuesto (precios) cuando:}
\begin{opts}
    \item La curva de daño marginal es muy empinada (riesgo de umbral catastrófico) respecto a la de costo de abatimiento.
    \item Los costos marginales de abatimiento son completamente conocidos con certeza antes de fijar la regulación.
    \item La curva de daño marginal ambiental es relativamente plana y la de costo marginal de abatimiento es muy empinada.
    \item Existen únicamente dos empresas reguladas con idénticas funciones de costos marginales constantes.
\end{opts}

\q{Cuando una empresa aguas arriba vierte efluentes que elevan los costos operativos de una empresa aguas abajo, la fusión de ambas bajo una sola propiedad logra la eficiencia porque:}
\begin{opts}
    \item La fusión elimina el poder de mercado en el producto final y garantiza un precio igual al costo marginal.
    \item La empresa integrada maximiza el beneficio conjunto, internalizando el daño externo en su decisión de producción.
    \item El Estado subsidia automáticamente las inversiones de tratamiento ambiental de las empresas verticalmente integradas.
    \item Las externalidades negativas desaparecen físicamente al compartirse la misma razón social tributaria.
\end{opts}

\q{En análisis económico del derecho ambiental, una regla de responsabilidad objetiva (strict liability), que obliga al contaminador a compensar todo daño causado independientemente de su nivel de precaución:}
\begin{opts}
    \item Exime al contaminador de pagar indemnizaciones siempre que demuestre haber cumplido el estándar legal mínimo.
    \item Reduce el incentivo de la empresa a invertir en tecnologías limpias porque la multa es independiente del daño.
    \item Induce al contaminador a elegir tanto el nivel óptimo de precaución como el nivel eficiente de actividad.
    \item Requiere que los tribunales conozcan con precisión el nivel de diligencia óptimo antes de dictar sentencia.
\end{opts}

\q{Los Mecanismos de Retribución por Servicios Ecosistémicos (MERESE) impulsados y regulados por SUNASS en las tarifas de las EPS constituyen un instrumento económico destinado a:}
\begin{opts}
    \item Financiar subsidios cruzados para el consumo industrial de agua potable en zonas costeras urbanas.
    \item Reemplazar el tratamiento químico en las plantas potabilizadoras por el libre acceso a los acuíferos.
    \item Subastar derechos perpetuos de propiedad privada sobre los ríos y lagunas a favor de operadores privados.
    \item Internalizar externalidades hidrológicas retribuyendo la conservación de cuencas altas.
\end{opts}

\q{En la regulación sanitaria peruana supervisada por SUNASS, el cobro de un 'pago adicional por exceso de concentración' a usuarios no domésticos que superan los Valores Máximos Admisibles (VMA) en sus descargas al alcantarillado opera económicamente como:}
\begin{opts}
    \item Un impuesto pigouviano que internaliza el mayor costo de operación y daño causado a redes y PTAR.
    \item Un precio de Lindahl voluntario basado en la revelación honesta de preferencias de cada industria.
    \item Un subsidio cruzado regresivo que transfiere recursos de los usuarios domésticos hacia las industrias.
    \item Una cuota de comando y control que prohíbe cualquier descarga industrial independientemente de su concentración.
\end{opts}

\q{El vertimiento de aguas residuales urbanas sin tratamiento a ríos o al mar constituye una falla de mercado clásica que justifica la inversión y regulación en Plantas de Tratamiento de Aguas Residuales (PTAR) porque:}
\begin{opts}
    \item El agua residual tratada es un bien privado puro con alta rentabilidad comercial inmediata sin regulación.
    \item Los usuarios individuales no internalizan el daño ambiental y sanitario impuesto a las poblaciones aguas abajo.
    \item Las PTAR presentan deseconomías de escala pronunciadas que favorecen la competencia entre múltiples plantas pequeñas.
    \item El costo marginal de construir una PTAR es siempre igual a cero una vez instalada la red de agua potable.
\end{opts}

\q{Desde la microeconomía del bienestar, la expansión de conexiones domiciliarias de agua potable y alcantarillado hacia hogares de bajos ingresos genera una externalidad positiva en salud pública porque:}
\begin{opts}
    \item Convierte el agua potable en un bien público puro no rival y no excluible en el consumo domiciliario.
    \item Disminuye el costo marginal privado de producción de la empresa prestadora hasta hacerlo estrictamente negativo.
    \item Reduce el contagio comunitario de enfermedades hídricas, elevando $BMS$ sobre $BMP$.
    \item Elimina la necesidad de instalar micromedidores al volver inelástica la demanda residencial de agua.
\end{opts}

\q{En políticas de conservación de recursos naturales (agua o energía), el 'efecto rebote' (o paradoja de Jevons) describe el fenómeno mediante el cual:}
\begin{opts}
    \item Un aumento en el impuesto pigouviano provoca una reducción más que proporcional en el consumo del recurso.
    \item La instalación de medidores inteligentes elimina por completo las pérdidas comerciales y técnicas en la red.
    \item Los consumidores reducen voluntariamente su demanda cuando perciben que el recurso común está próximo a agotarse.
    \item Una mejora en eficiencia abarata el uso efectivo y puede elevar el consumo total.
\end{opts}

\q{Si una empresa puede reducir sus emisiones contaminantes tanto disminuyendo su producción como instalando filtros de abatimiento (end-of-pipe), un impuesto directo sobre las emisiones resulta superior a un impuesto sobre el producto final porque:}
\begin{opts}
    \item El impuesto al producto solo reduce escala y no incentiva instalar abatimiento.
    \item El impuesto sobre el producto siempre recauda menos ingresos fiscales que el impuesto directo sobre emisiones.
    \item El impuesto a las emisiones no requiere ningún sistema de monitoreo ni medición de los efluentes vertidos.
    \item El impuesto sobre el producto eleva los costos marginales de abatimiento por encima del daño marginal social.
\end{opts}

\q{En una carretera urbana de libre acceso sujeta a congestión, el equilibrio sin peajes presenta exceso de tráfico porque cada conductor decide ingresar comparando su beneficio marginal con:}
\begin{opts}
    \item El costo marginal social de congestión que su vehículo impone sobre todos los demás conductores.
    \item El costo medio privado del viaje, ignorando la demora impuesta a otros usuarios.
    \item El costo fijo de construcción de la infraestructura vial dividido entre la población total de la ciudad.
    \item El costo marginal de mantenimiento del pavimento cuando el flujo vehicular tiende a cero.
\end{opts}

\q{En la evaluación de externalidades ambientales intergeneracionales (como el cambio climático o infraestructura hídrica de largo plazo), la regla de Ramsey expresa la tasa social de descuento como $r = \rho + \eta g$, donde:}
\begin{opts}
    \item $\rho$ es el impuesto pigouviano óptimo, $\eta$ la elasticidad-precio de la demanda y $g$ el crecimiento poblacional.
    \item $\rho$ es la tasa de inflación esperada, $\eta$ el índice de Lerner del sector y $g$ la tasa de depreciación del capital.
    \item $\rho$ es impaciencia pura, $\eta$ elasticidad de utilidad marginal y $g$ crecimiento del consumo.
    \item $\rho$ es la productividad marginal del capital, $\eta$ la aversión absoluta al riesgo y $g$ el déficit fiscal.
\end{opts}

\q{Al comparar el principio 'quien contamina paga' (impuesto pigouviano sin compensación directa a las víctimas) frente a compensar a las víctimas proporcionalmente al daño sufrido, la teoría económica advierte que compensar directamente a las víctimas puede:}
\begin{opts}
    \item Reducir los incentivos de la empresa contaminadora para invertir en tecnologías de abatimiento.
    \item Asegurar que los costos de transacción se vuelvan nulos en cualquier conflicto ambiental multilateral.
    \item Aumentar la recaudación neta del gobierno al eliminar los costos administrativos de fiscalización ambiental.
    \item Generar riesgo moral en las víctimas, reduciendo su incentivo a autoprotegerse.
\end{opts}

\q{En la teoría de los bienes club de Buchanan (bienes excluibles pero no rivales hasta cierto umbral de congestión), la membresía óptima $N^*$ del club para un tamaño dado de la instalación se alcanza cuando:}
\begin{opts}
    \item El ahorro en costo por socio al sumar un miembro iguala su costo de congestión.
    \item El número de miembros es igual a uno para garantizar la ausencia absoluta de congestión en el uso del servicio.
    \item El costo medio por miembro alcanza su valor máximo y la disposición marginal a pagar cae exactamente a cero.
    \item Todos los ciudadanos del país ingresan gratuitamente al club financiando la infraestructura con impuestos generales.
\end{opts}

\q{En el modelo de provisión privada voluntaria de un bien público puro (Bergstrom, Blume y Varian), si el Estado grava con un impuesto de suma alzada a los contribuyentes activos para proveer directamente el bien público, el teorema de neutralidad de Warr implica que:}
\begin{opts}
    \item La provisión total del bien público aumenta exactamente en el mismo monto que el gasto público adicional.
    \item Los privados reducen sus aportes sol por sol, generando crowding-out completo.
    \item El problema del polizón (free-rider) desaparece de inmediato y se alcanza la condición eficiente de Samuelson.
    \item Los individuos no contribuyentes incrementan voluntariamente sus donaciones al observar el esfuerzo estatal.
\end{opts}

\q{Cuando un bien público puro debe financiarse mediante impuestos distorsionadores (en lugar de impuestos de suma alzada), la regla modificada de Samuelson (Atkinson y Stern) establece que la suma de las $TMS_i$ debe igualar:}
\begin{opts}
    \item Estrictamente a cero para compensar la pérdida de eficiencia provocada por la tributación distorsionadora.
    \item Al costo marginal de producción dividido entre la elasticidad-precio de la oferta del bien público.
    \item Al costo marginal de producción multiplicado por el costo marginal de los fondos públicos ($MCPF > 1$).
    \item Al costo medio total del bien público menos la recaudación obtenida por impuestos pigouvianos.
\end{opts}

\q{La razón económica fundamental por la cual la demanda social de un bien público puro se construye mediante la suma vertical de las demandas individuales (y no horizontal, como en bienes privados) es que:}
\begin{opts}
    \item Todos los consumidores pagan exactamente el mismo precio unitario pero consumen cantidades distintas del bien.
    \item Los consumidores pueden excluirse mutuamente del disfrute del bien mediante contratos privados bilaterales.
    \item El costo marginal de producir una unidad adicional del bien público es siempre decreciente en el largo plazo.
    \item Al ser no rival, todos consumen la misma cantidad $G$ y se suman sus valoraciones.
\end{opts}

\q{En el modelo de bienes públicos locales de Tiebout (1956), la movilidad residencial sin costos entre múltiples jurisdicciones con distintas combinaciones de impuestos y servicios ('votar con los pies') permite:}
\begin{opts}
    \item Inducir a revelar preferencias por bienes públicos locales al elegir residencia.
    \item Eliminar la rivalidad en el consumo de los bienes privados producidos dentro de cada municipalidad.
    \item Garantizar la igualación absoluta de los ingresos per cápita entre todas las jurisdicciones locales.
    \item Resolver las externalidades ambientales globales sin necesidad de acuerdos internacionales vinculantes.
\end{opts}

\q{Aunque el equilibrio de Lindahl cumple la condición de Samuelson mediante precios personalizados $\tau_i = TMS_i$, su principal falla práctica para implementarse de forma descentralizada consiste en que:}
\begin{opts}
    \item Genera un déficit presupuestario permanente porque la suma de precios personalizados es inferior al costo marginal.
    \item Cada agente tiene incentivos a subdeclarar su verdadera valoración para pagar un precio personalizado menor.
    \item Asigna cantidades distintas del bien público a cada ciudadano según su nivel de ingreso disponible.
    \item Requiere que todos los individuos posean exactamente la misma función de utilidad y la misma renta.
\end{opts}

\q{En el mecanismo de revelación de preferencias de Vickrey-Clarke-Groves (VCG) para decidir la provisión de un bien público, decir la verdad es una estrategia débilmente dominante porque:}
\begin{opts}
    \item El gobierno distribuye íntegramente los impuestos de pivote recaudados entre los mismos votantes participantes.
    \item Ningún ciudadano paga impuestos si el proyecto del bien público resulta finalmente aprobado por mayoría.
    \item Solo paga si su voto es decisivo (pivote), cubriendo el daño impuesto a otros.
    \item Se prohíbe votar a los individuos cuya valoración neta por el bien público sea estrictamente negativa.
\end{opts}

\q{El Teorema de Gibbard-Satterthwaite en teoría de elección social establece que, con al menos tres alternativas y preferencias no restringidas, toda regla determinística de elección colectiva que sea a prueba de manipulación estratégica (strategy-proof) debe ser:}
\begin{opts}
    \item Unánime y necesariamente financiada mediante impuestos de pivote de Clarke-Groves.
    \item Basada en la regla de mayoría simple por pares para evitar ciclos de Condorcet.
    \item Invariante frente a la eliminación de alternativas irrelevantes únicamente con dos votantes.
    \item Dictatorial, es decir, reflejar siempre la alternativa preferida de un único individuo.
\end{opts}

\q{En un acuífero subterráneo de libre acceso (recurso común), la sobreexplotación conocida como la 'Tragedia de los Comunes' ocurre porque cada agricultor o usuario extrae agua hasta que su costo marginal de extracción iguala:}
\begin{opts}
    \item El producto medio ($PMe$), ignorando la caída de rendimiento para los demás.
    \item El producto marginal social del recurso ($PMg$), ignorando las economías de escala en la perforación.
    \item El precio de monopolio del agua embotellada en el mercado minorista más cercano al acuífero.
    \item La suma vertical de las disposiciones marginales a pagar de todos los usuarios del acuífero.
\end{opts}

\q{Los estudios empíricos y teóricos de Elinor Ostrom (Premio Nobel 2009) sobre sistemas de riego y recursos comunes demostraron que la Tragedia de los Comunes:}
\begin{opts}
    \item Solo puede resolverse mediante la estatización centralizada total o la privatización absoluta del recurso.
    \item Puede evitarse con autogestión comunitaria, reglas claras, monitoreo y sanciones.
    \item Es matemáticamente inevitable siempre que existan más de dos usuarios explotando un recurso no excluible.
    \item Ocurre únicamente cuando el recurso posee simultáneamente propiedades de no rivalidad y alta exclusión.
\end{opts}

\q{El otorgamiento de licencias con volúmenes máximos de extracción y tarifas por uso de agua subterránea en acuíferos sobreexplotados (como en valles costeros peruanos) corrige la falla del recurso común porque:}
\begin{opts}
    \item Permite que cada usuario extraiga agua ilimitadamente sin pagar por el descenso del nivel freático.
    \item Transforma físicamente el agua subterránea en un bien público puro con costo marginal nulo de extracción.
    \item Introduce exclusión y un precio sombra que internaliza el agotamiento compartido.
    \item Elimina la rivalidad entre los distintos agricultores que bombean agua del mismo acuífero.
\end{opts}

\q{Si un bien es no rival hasta su capacidad máxima pero resulta técnicamente excluible (por ejemplo, una red de alcantarillado pluvial o un puente sin congestión), cobrar un precio positivo por su uso cuando está descongestionado genera:}
\begin{opts}
    \item Una asignación eficiente de Pareto al cubrir automáticamente los costos fijos de construcción sin distorsión.
    \item Una sobreexplotación del recurso idéntica a la Tragedia de los Comunes en acuíferos abiertos.
    \item Un incentivo excesivo al problema del polizón (free-rider) entre los usuarios que pagan el peaje.
    \item Una pérdida de eficiencia social porque se excluye a usuarios cuya valoración supera el costo marginal nulo.
\end{opts}

\q{En economía del sector público, la distinción entre 'provisión pública' y 'producción pública' de un bien público (o infraestructura sanitaria) implica que:}
\begin{opts}
    \item El Estado puede financiar la provisión pero delegar la producción a privados.
    \item Todo bien financiado con recursos públicos debe ser construido y operado exclusivamente por empresas estatales.
    \item La producción pública garantiza siempre costos menores que la contratación privada en licitaciones competitivas.
    \item La provisión privada voluntaria es idéntica a la producción estatal bajo concesiones de largo plazo.
\end{opts}

\q{El concepto de 'bienes meritorios' (merit goods) introducido por Richard Musgrave justifica la intervención estatal (por ejemplo, subsidiar el consumo mínimo vital de agua potable o educación básica) con base en:}
\begin{opts}
    \item La no rivalidad y no exclusión técnica absoluta que impide cobrar cualquier tarifa a los usuarios.
    \item Que los individuos subvaloran el bien por miopía, más allá de externalidades.
    \item La existencia de rendimientos decrecientes a escala que impiden la rentabilidad de las empresas privadas.
    \item La aplicación estricta de la soberanía del consumidor sin ningún componente de paternalismo estatal.
\end{opts}

\q{Cuando una comunidad decide por votación de mayoría simple el nivel de gasto en un bien público unidimensional, el Teorema del Votante Mediano garantiza que la alternativa preferida por el votante mediano vence a cualquier otra si:}
\begin{opts}
    \item Existen al menos tres dimensiones independientes de gasto público votadas simultáneamente.
    \item Todos los ciudadanos pagan exactamente el mismo impuesto de Lindahl proporcional a su valoración.
    \item Todos los votantes tienen preferencias de pico único sobre el nivel de gasto.
    \item Las preferencias de los votantes son multimodales y generan ciclos intransitivos de Condorcet.
\end{opts}

\q{El Teorema de Imposibilidad de Arrow demuestra que, con tres o más alternativas sociales, no existe ninguna función de bienestar social que agregue preferencias individuales racionales cumpliendo simultáneamente:}
\begin{opts}
    \item Preferencias cuasilineales, cero costos de transacción, información completa y pagos de transferencia.
    \item Equilibrio presupuestario, condición de Samuelson, simetría de costos y ausencia de externalidades.
    \item Dominio restringido a pico único, regla de mayoría simple, anonimato y neutralidad de alternativas.
    \item Dominio irrestricto, eficiencia de Pareto, independencia de alternativas irrelevantes y no dictadura.
\end{opts}

\q{Aunque ambos involucran interdependencias en el consumo, una externalidad de red directa (por ejemplo, telefonía o mensajería) se diferencia de un bien público puro porque:}
\begin{opts}
    \item El acceso es excluible y su valor crece con el número de usuarios conectados.
    \item El bien público puro siempre tiene un valor creciente con el número de usuarios y nunca sufre congestión.
    \item La externalidad de red solo puede ser provista por el Estado mediante impuestos generales obligatorios.
    \item El bien público puro es estrictamente rival en el consumo pero excluible mediante tarifas en dos partes.
\end{opts}

\q{En contraste con la Tragedia de los Comunes (sobreexplotación por falta de exclusión), la 'Tragedia de los Anticomunes' formulada por Michael Heller ocurre cuando:}
\begin{opts}
    \item El Estado elimina los derechos de propiedad privada en una cuenca hídrica causando su agotamiento inmediato.
    \item Exceso de derechos de exclusión fragmentados que provoca subutilización del recurso.
    \item Un monopolista único posee todos los derechos de propiedad sobre una patente y cobra precios discriminatorios.
    \item Los consumidores sobreexplotan un bien público puro debido a la ausencia de precios de Lindahl.
\end{opts}

\q{En el modelo de selección adversa de Akerlof (mercado de 'limones'), el mecanismo de desenredo (unraveling) que puede llevar al colapso del intercambio de bienes de alta calidad ocurre porque:}
\begin{opts}
    \item Los compradores observan la calidad real pero se coluden para ofrecer un precio inferior al costo marginal.
    \item Los vendedores de mala calidad se retiran primero del mercado al no poder cubrir sus costos fijos hundidos.
    \item El precio promedio expulsa a los bienes buenos, bajando aún más la calidad media.
    \item El gobierno fija un precio mínimo obligatorio que impide la venta de vehículos usados de baja calidad.
\end{opts}

\q{Según el modelo de Stiglitz y Weiss (1981), ante un exceso de demanda de préstamos los bancos pueden optar por racionar el crédito en lugar de subir la tasa de interés porque una tasa más alta:}
\begin{opts}
    \item Aumenta la probabilidad de repago de todos los proyectos al incentivar mayor esfuerzo gerencial.
    \item Está prohibida universalmente por el Primer Teorema del Bienestar en mercados financieros competitivos.
    \item Elimina la asimetría de información entre prestamistas y prestatarios al obligar a revelar el riesgo real.
    \item Atrae deudores riesgosos (selección adversa) e induce proyectos más arriesgados.
\end{opts}

\q{En el modelo de señalización educativa de Spence, la condición de 'cruce único' (single-crossing) —indispensable para que exista un equilibrio separador— requiere que:}
\begin{opts}
    \item El costo marginal de adquirir una unidad adicional de señal sea estrictamente menor para el tipo de alta productividad.
    \item La educación incremente necesariamente la productividad física marginal de todos los trabajadores en igual magnitud.
    \item Los empleadores conozcan ex ante la productividad real de cada trabajador antes de observar su nivel educativo.
    \item Los trabajadores de baja productividad tengan una mayor disposición a pagar por educación que los de alta productividad.
\end{opts}

\q{En el modelo puro de señalización de Spence donde la educación no aumenta la productividad laboral sino que solo actúa como señal, un equilibrio separador puede resultar socialmente ineficiente respecto a un equilibrio agrupador porque:}
\begin{opts}
    \item Los empleadores obtienen beneficios extraordinarios positivos pagando salarios inferiores a la productividad media.
    \item La señal consume recursos reales sin elevar el producto, pudiendo reducir bienestar.
    \item Los trabajadores de baja productividad se ven obligados a invertir el doble en educación que los de alta productividad.
    \item El equilibrio separador viola necesariamente la restricción de compatibilidad de incentivos de los trabajadores.
\end{opts}

\q{En el modelo competitivo de cribado (screening) en seguros de Rothschild y Stiglitz (1976), nunca puede existir un equilibrio agrupador (pooling) donde altos y bajos riesgos compren la misma póliza porque:}
\begin{opts}
    \item Los individuos de alto riesgo siempre prefieren quedarse sin seguro antes que pagar una prima actuarialmente justa.
    \item Las compañías de seguros son aversas al riesgo y no pueden diversificar pólizas de distintos perfiles.
    \item Un rival puede ofrecer menor cobertura y menor prima atrayendo solo bajos riesgos.
    \item La regulación estatal prohíbe ofrecer contratos de seguro con cobertura completa a los clientes de alto riesgo.
\end{opts}

\q{Cuando existe un equilibrio separador competitivo en el modelo de seguros de Rothschild y Stiglitz, las características de los dos contratos ofrecidos en equilibrio son:}
\begin{opts}
    \item Cobertura parcial para ambos tipos de asegurados con una prima única igual al riesgo promedio de la población.
    \item Cobertura completa para el tipo de bajo riesgo y cobertura incompleta (con deducible) para el tipo de alto riesgo.
    \item Cobertura completa para ambos tipos pagando exactamente la misma prima actuarialmente justa del tipo bajo.
    \item Cobertura completa para el tipo de alto riesgo y cobertura incompleta (con deducible) para el tipo de bajo riesgo.
\end{opts}

\q{En el modelo de Rothschild y Stiglitz, ¿en qué circunstancia puede no existir siquiera un equilibrio separador de Nash en estrategias puras?}
\begin{opts}
    \item Cuando hay muchos bajos riesgos y un contrato agrupador atrae a ambos tipos.
    \item Cuando la proporción de individuos de alto riesgo es igual al 100\% de la población asegurada.
    \item Cuando las aseguradoras observan sin costo el nivel de riesgo individual antes de firmar el contrato.
    \item Cuando todos los consumidores son estrictamente neutrales al riesgo frente a la pérdida patrimonial.
\end{opts}

\q{En un mercado de seguros con riesgo moral (acción de prevención no observable por la aseguradora), el contrato óptimo de Second-Best incluye deducibles o copagos porque:}
\begin{opts}
    \item Permiten a la aseguradora obtener beneficios monopólicos aun bajo libre entrada en el mercado asegurador.
    \item El asegurado asume parte del daño, balanceando cobertura e incentivo a prevenir.
    \item Logran exactamente la misma asignación eficiente de First-Best que existiría si la prevención fuera observable.
    \item Eliminan la aversión al riesgo de los asegurados al transformar su función de utilidad en lineal.
\end{opts}

\q{En el modelo Principal-Agente con riesgo moral (esfuerzo no observable), si el agente es neutral al riesgo y no tiene restricciones de liquidez, el principal puede alcanzar el óptimo de First-Best mediante:}
\begin{opts}
    \item Un salario fijo independiente del resultado que asegure completamente el ingreso del agente.
    \item Una participación accionaria del 50\% para cada parte independientemente de la varianza del resultado.
    \item Un contrato de franquicia con renta fija donde el agente retiene todo el ingreso.
    \item La supervisión directa continua del esfuerzo del agente mediante un contrato de costo más margen.
\end{opts}

\q{En el modelo Principal-Agente con esfuerzo no observable donde el principal es neutral al riesgo y el agente es averso al riesgo, el contrato óptimo de Second-Best:}
\begin{opts}
    \item Transfiere el 100\% del riesgo de producción al agente para maximizar su esfuerzo sin costo de agencia.
    \item Asigna un salario completamente fijo al agente para evitar pagarle cualquier prima por riesgo.
    \item Garantiza que la restricción de compatibilidad de incentivos siempre sea redundante y no vinculante.
    \item Vincula parcialmente el pago al resultado, balanceando incentivos y riesgo.
\end{opts}

\q{El Principio de Informatividad de Holmström (1979) en teoría de contratos establece que un esquema de remuneración óptimo para el agente debe depender de una variable observable adicional $z$ si y solo si:}
\begin{opts}
    \item Aporta información marginal sobre el esfuerzo más allá del resultado observado $y$.
    \item La variable $z$ no tiene correlación alguna con el esfuerzo del agente pero sí con el riesgo macroeconómico.
    \item El agente puede manipular directamente el valor de la variable $z$ sin incrementar su costo de esfuerzo.
    \item El principal es más averso al riesgo que el agente y necesita diversificar sus ingresos financieros.
\end{opts}

\q{Según el modelo de agencia multitarea de Holmström y Milgrom (1991), cuando un agente (o una EPS regulada) realiza simultáneamente una tarea fácil de medir (recortar gastos) y otra difícil de medir (calidad del agua), otorgar incentivos monetarios muy fuertes solo a la tarea medible:}
\begin{opts}
    \item Mejora automáticamente el desempeño en ambas tareas gracias a las complementariedades gerenciales.
    \item Desvía esfuerzo hacia la tarea medible dañando calidad; convienen incentivos moderados.
    \item Elimina por completo el problema de riesgo moral sin necesidad de fijar estándares mínimos de calidad.
    \item Alcanza el First-Best siempre que el agente sea estrictamente averso al riesgo.
\end{opts}

\q{En el modelo de riesgo moral laboral de Shapiro y Stiglitz (1984), las empresas pagan voluntariamente un 'salario de eficiencia' por encima del salario que vacía el mercado con el fin de:}
\begin{opts}
    \item Señalizar a los consumidores finales que la empresa opera bajo competencia perfecta con beneficios nulos.
    \item Eliminar por completo el desempleo involuntario en el equilibrio de largo plazo del mercado de trabajo.
    \item Crear un costo de despido (desempleo involuntario) que disuada el escaqueo.
    \item Aumentar la rotación laboral voluntaria de los trabajadores de mayor productividad hacia otras firmas.
\end{opts}

\q{En el modelo de regulación de un monopolio con información privada sobre sus costos (Baron y Myerson, 1982), el contrato óptimo del regulador se caracteriza por:}
\begin{opts}
    \item Extraer toda la renta de la empresa eficiente y producir la cantidad de First-Best en ambos tipos de costos.
    \item Fijar un precio único igual al costo marginal esperado sin ofrecer un menú diferenciado de contratos.
    \item Conceder una renta informacional a la firma ineficiente y sobreproducir cuando la firma es de bajo costo.
    \item Dar renta informacional al tipo eficiente y reducir la producción del ineficiente.
\end{opts}

\q{En el modelo de regulación de Laffont y Tirole (1986), que combina selección adversa (eficiencia intrínseca oculta) y riesgo moral (esfuerzo reductor de costos no observable), el regulador óptimo ofrece un menú de contratos donde:}
\begin{opts}
    \item El tipo eficiente elige contrato de alto poder y el ineficiente costo compartido.
    \item Todas las empresas reciben un contrato de costo-plus puro (bajo poder) para eliminar cualquier renta informacional.
    \item Todas las empresas reciben un precio tope idéntico sin transferencia alguna, logrando cero rentas informacionales.
    \item El regulador renuncia a observar los costos contables realizados para evitar cualquier distorsión del esfuerzo.
\end{opts}

\q{El Principio de Revelación en diseño de mecanismos establece que, para encontrar el contrato o mecanismo regulatorio óptimo bajo información asimétrica, el diseñador puede restringirse sin pérdida de generalidad a:}
\begin{opts}
    \item Mecanismos de subasta inglesa donde todos los agentes conocen los tipos privados de sus competidores.
    \item Mecanismos directos donde revelar verazmente el tipo es compatible con incentivos.
    \item Contratos lineales de precio único que asignan exactamente la misma cantidad a todos los tipos de agentes.
    \item Esquemas de fiscalización aleatoria que eliminan la necesidad de satisfacer la restricción de participación.
\end{opts}

\q{De acuerdo con la teoría económica de la regulación (Stigler-Peltzman) y el enfoque de agencia (Laffont-Tirole), el riesgo de 'captura regulatoria' surge principalmente porque:}
\begin{opts}
    \item Los consumidores están altamente concentrados y poseen mayor información técnica que las empresas reguladas.
    \item Los organismos reguladores fijan tarifas iguales al costo marginal sin requerir información contable de la firma.
    \item La industria tiene intereses concentrados para influir frente a usuarios dispersos.
    \item La competencia perfecta en monopolios naturales elimina los beneficios extraordinarios susceptibles de captura.
\end{opts}

\q{Cuando una Empresa Prestadora de Servicios de Saneamiento (EPS) presenta su Plan Maestro Optimizado (PMO) ante SUNASS, existe un problema de selección adversa porque la EPS puede:}
\begin{opts}
    \item Reducir deliberadamente la calidad del agua después de haberse aprobado la fórmula tarifaria quinquenal.
    \item Transferir automáticamente todas sus pérdidas técnicas de agua a una aseguradora privada competitiva.
    \item Competir agresivamente en precios con otra EPS vecina para capturar a los usuarios de mayor consumo.
    \item Sobredeclarar costos operativos o inversiones aprovechando su ventaja informativa.
\end{opts}

\q{Bajo un esquema de regulación por precio tope (price cap) en servicios de saneamiento, el principal riesgo moral post-contractual que debe monitorear SUNASS mediante metas de gestión consiste en que la EPS:}
\begin{opts}
    \item Subinvertir en mantenimiento y degradar la calidad para recortar costos operativos.
    \item Sobreinvierta innecesariamente en capital físico para inflar su base de activos (efecto Averch-Johnson).
    \item Revele voluntariamente menores costos históricos de los reales durante el proceso de revisión tarifaria.
    \item Elimine la micromedición residencial para reducir la facturación mensual de los usuarios domésticos.
\end{opts}

\q{En mercados de bienes durables con calidad oculta, el ofrecimiento de una garantía amplia por parte del vendedor actúa como una señal creíble de alta calidad (análogo a Spence) porque:}
\begin{opts}
    \item Las garantías eliminan por completo la posibilidad de riesgo moral por mal uso del comprador.
    \item Cumplir la garantía cuesta mucho menos para un producto bueno que para uno malo.
    \item Solo las empresas con productos de baja calidad están dispuestas a ofrecer garantías extendidas gratuitas.
    \item La ley obliga a todos los fabricantes a cobrar exactamente el mismo precio independientemente de la garantía.
\end{opts}

\q{En concesiones de infraestructura sanitaria con inversiones específicas y contratos incompletos, el problema de retención u oportunismo ex post (hold-up) provoca ex ante:}
\begin{opts}
    \item Una sobreinversión sistemática en activos específicos debido a la seguridad jurídica absoluta del contrato.
    \item La eliminación inmediata de los costos hundidos y la transformación del sector en un mercado desafiable.
    \item Subinversión en activos específicos por temor a expropiación ex post de rentas.
    \item Un aumento de la eficiencia asignativa al permitir renegociar libremente las tarifas cada trimestre.
\end{opts}

\subsection*{Parte II. Aplicaciones prácticas y ejercicios cuantitativos (154 a 180)}

\q{Un mercado competitivo presenta demanda $P = 100 - Q$, costo marginal privado $CMP = 10 + Q$ y costo marginal externo $CME = 0.5Q$. La cantidad de mercado $Q_M$, el óptimo social $Q_{soc}^*$ y el impuesto pigouviano óptimo $t^*$ son:}
\begin{opts}
    \item $Q_M = 50$, $Q_{soc}^* = 40$ y $t^* = 20.0$
    \item $Q_M = 45$, $Q_{soc}^* = 36$ y $t^* = 22.5$
    \item $Q_M = 45$, $Q_{soc}^* = 30$ y $t^* = 15.0$
    \item $Q_M = 45$, $Q_{soc}^* = 36$ y $t^* = 18.0$
\end{opts}

\q{Una industria competitiva tiene demanda $P = 80 - Q$, costo marginal privado constante $CMP = 20$ y genera una externalidad negativa constante $CME = 12$ por unidad. La Pérdida de Eficiencia Social ($PIE$) en el equilibrio sin regulación es:}
\begin{opts}
    \item $PIE = 72$
    \item $PIE = 36$
    \item $PIE = 108$
    \item $PIE = 144$
\end{opts}

\q{En un mercado competitivo con demanda $P = 120 - 2Q$, $CMP = 20 + Q$ y costo marginal externo $CME = Q$, el Estado introduce el impuesto pigouviano óptimo $t^*$. La recaudación tributaria total ($R = t^* Q_{soc}^*$) asciende a:}
\begin{opts}
    \item $R = 400$
    \item $R = 625$
    \item $R = 500$
    \item $R = 750$
\end{opts}

\q{El consumo de conexiones de saneamiento presenta beneficio marginal privado $BMP = 60 - Q$, beneficio marginal externo constante $BME = 18$ y costo marginal $CMg = 12 + Q$. La cantidad de mercado $Q_M$, el óptimo social $Q_{soc}^*$ y el subsidio óptimo $s^*$ son:}
\begin{opts}
    \item $Q_M = 24$, $Q_{soc}^* = 30$ y $s^* = 12$
    \item $Q_M = 24$, $Q_{soc}^* = 36$ y $s^* = 18$
    \item $Q_M = 24$, $Q_{soc}^* = 33$ y $s^* = 18$
    \item $Q_M = 30$, $Q_{soc}^* = 39$ y $s^* = 18$
\end{opts}

\q{Un bien con externalidad positiva en consumo tiene $BMP = 50 - Q$, $BME = 20$ constante y $CMg = 10 + Q$. La Pérdida de Eficiencia Social ($PIE$) ocasionada por el subconsumo en el equilibrio de mercado libre es:}
\begin{opts}
    \item $PIE = 250$
    \item $PIE = 150$
    \item $PIE = 200$
    \item $PIE = 100$
\end{opts}

\q{Un monopolista enfrenta la demanda $P = 100 - Q$, tiene costo marginal privado $CMP = 20$ y genera una externalidad negativa constante $CME = 10$ por unidad. Al comparar la cantidad monopólica $Q_M$ con el óptimo social $Q_{soc}^*$, se obtiene:}
\begin{opts}
    \item $Q_M = 40$ y $Q_{soc}^* = 70$, por lo que existe subproducción pese a la externalidad.
    \item $Q_M = 70$ y $Q_{soc}^* = 40$, por lo que se requiere un impuesto pigouviano de $t = 10$.
    \item $Q_M = 40$ y $Q_{soc}^* = 35$, por lo que existe sobreproducción neta de $5$ unidades.
    \item $Q_M = 50$ y $Q_{soc}^* = 50$, compensándose exactamente ambas fallas de mercado.
\end{opts}

\q{Una fábrica vierte efluentes que causan un daño de $S/\ 600$ a una piscigranja aguas abajo; instalar un filtro que elimina totalmente el daño cuesta $S/\ 350$ a la fábrica. Si los derechos de propiedad pertenecen a la fábrica y no hay costos de transacción, la negociación coasiana logra el filtro mediante un pago $T$ de la piscigranja en el rango:}
\begin{opts}
    \item $S/\ 0 \le T \le S/\ 250$
    \item $S/\ 350 \le T \le S/\ 600$
    \item $S/\ 250 \le T \le S/\ 350$
    \item $S/\ 600 \le T \le S/\ 950$
\end{opts}

\q{Dos empresas deben reducir sus emisiones en $40$ toneladas en total ($a_1 + a_2 = 40$) bajo un sistema de permisos negociables. Sus costos marginales de abatimiento son $CMgA_1 = 2a_1$ y $CMgA_2 = 6a_2$. En el equilibrio eficiente, las reducciones $(a_1^*, a_2^*)$ y el precio del permiso $p^*$ son:}
\begin{opts}
    \item $a_1^* = 20$, $a_2^* = 20$ y $p^* = 40$
    \item $a_1^* = 24$, $a_2^* = 16$ y $p^* = 48$
    \item $a_1^* = 30$, $a_2^* = 10$ y $p^* = 60$
    \item $a_1^* = 32$, $a_2^* = 8$ y $p^* = 64$
\end{opts}

\q{Dos firmas con costos totales de abatimiento $CTA_1 = a_1^2$ y $CTA_2 = 3a_2^2$ deben lograr una reducción conjunta de $40$ toneladas. Comparando un sistema de permisos negociables frente a un estándar uniforme ($a_1 = a_2 = 20$), el ahorro total de costos sociales es:}
\begin{opts}
    \item $200$ unidades monetarias
    \item $300$ unidades monetarias
    \item $600$ unidades monetarias
    \item $400$ unidades monetarias
\end{opts}

\q{Una planta industrial obtiene un beneficio marginal $BMg = 90 - 3Q$ al descargar $Q$ unidades de carga orgánica al alcantarillado, generando un daño marginal en la PTAR de $DMg = 15 + 2Q$. El nivel eficiente de descarga $Q^*$ y el cargo unitario óptimo $t^*$ son:}
\begin{opts}
    \item $Q^* = 15$ y $t^* = 45$
    \item $Q^* = 18$ y $t^* = 36$
    \item $Q^* = 20$ y $t^* = 55$
    \item $Q^* = 25$ y $t^* = 15$
\end{opts}

\q{Dos vecinos valoran un sistema de alumbrado público ($G$) según sus disposiciones marginales a pagar $DAP_1 = 40 - G$ y $DAP_2 = 60 - 2G$. Si el costo marginal de provisión es constante e igual a $CMg = 25$, la cantidad eficiente $G^*$ según la condición de Samuelson es:}
\begin{opts}
    \item $G^* = 15$ unidades
    \item $G^* = 25$ unidades
    \item $G^* = 20$ unidades
    \item $G^* = 30$ unidades
\end{opts}

\q{Una comunidad está conformada por $3$ hogares idénticos, cada uno con disposición marginal a pagar por protección contra inundaciones $DAP_i = 30 - 0.5G$. Si el costo marginal del proyecto es $CMg(G) = 30 + G$, el nivel socialmente óptimo $G^*$ es:}
\begin{opts}
    \item $G^* = 18$ unidades
    \item $G^* = 20$ unidades
    \item $G^* = 24$ unidades
    \item $G^* = 30$ unidades
\end{opts}

\q{Dos individuos tienen disposiciones marginales a pagar por un bien público puro de $DAP_1 = 50 - G$ y $DAP_2 = 40 - G$, cuyo costo marginal es $CMg = 20$. La provisión privada en Equilibrio de Nash ($G_{\text{priv}}$) y el óptimo de Samuelson ($G_{soc}^*$) son:}
\begin{opts}
    \item $G_{\text{priv}} = 20$ y $G_{soc}^* = 30$
    \item $G_{\text{priv}} = 25$ y $G_{soc}^* = 35$
    \item $G_{\text{priv}} = 30$ y $G_{soc}^* = 45$
    \item $G_{\text{priv}} = 30$ y $G_{soc}^* = 35$
\end{opts}

\q{Dos ciudadanos valoran un parque público con $DAP_1 = 40 - G$ y $DAP_2 = 60 - 2G$, siendo el costo marginal $CMg = 25$. En el Equilibrio de Lindahl, los precios personalizados $(\tau_1, \tau_2)$ por unidad de $G^*$ son:}
\begin{opts}
    \item $\tau_1 = 15.0$ y $\tau_2 = 10.0$
    \item $\tau_1 = 10.0$ y $\tau_2 = 15.0$
    \item $\tau_1 = 12.5$ y $\tau_2 = 12.5$
    \item $\tau_1 = 20.0$ y $\tau_2 = 5.0$
\end{opts}

\q{Tres vecinos deciden bajo un mecanismo de Clarke-Groves (VCG) si ejecutar un proyecto público. Sus beneficios netos (valoración menos cuota fija de costo) declarados son: Ana $= +80$, Bruno $= +40$ y Carlos $= -90$. Según las reglas VCG, el resultado y el impuesto de pivote que paga Ana son:}
\begin{opts}
    \item Se aprueba el proyecto y Ana paga un impuesto de pivote de $30$.
    \item Se aprueba el proyecto y Ana paga un impuesto de pivote de $50$.
    \item Se aprueba el proyecto y Ana paga un impuesto de pivote de $80$.
    \item Se rechaza el proyecto y ningún vecino paga impuesto de pivote.
\end{opts}

\q{Un acuífero de libre acceso produce un valor total de agua extraída $Y(N) = 100N - 2N^2$ en función del número de pozos $N$, donde operar cada pozo cuesta $c = 20$. El número de pozos en libre acceso ($N_{\text{libre}}$) y el número socialmente óptimo ($N_{soc}^*$) son:}
\begin{opts}
    \item $N_{\text{libre}} = 40$ y $N_{soc}^* = 30$
    \item $N_{\text{libre}} = 50$ y $N_{soc}^* = 25$
    \item $N_{\text{libre}} = 40$ y $N_{soc}^* = 20$
    \item $N_{\text{libre}} = 30$ y $N_{soc}^* = 15$
\end{opts}

\q{En un acuífero común con valor total de extracción $Y(N) = 100N - 2N^2$ y costo privado $c = 20$ por pozo, la autoridad hídrica desea cobrar una tarifa de licencia $t^*$ por pozo para inducir exactamente el número socialmente óptimo de pozos ($N_{soc}^* = 20$). Dicha tarifa debe ser:}
\begin{opts}
    \item $t^* = 20$ unidades monetarias
    \item $t^* = 30$ unidades monetarias
    \item $t^* = 60$ unidades monetarias
    \item $t^* = 40$ unidades monetarias
\end{opts}

\q{En una vía urbana congestionada, el costo medio privado por viaje es $CMe(x) = 10 + 0.5x$ (siendo $x$ el flujo vehicular) y la demanda inversa es $P = 70 - x$. El flujo socialmente óptimo $x_{soc}^*$ y el peaje óptimo de congestión $\tau^*$ son:}
\begin{opts}
    \item $x_{soc}^* = 30$ y $\tau^* = 15$
    \item $x_{soc}^* = 40$ y $\tau^* = 20$
    \item $x_{soc}^* = 30$ y $\tau^* = 25$
    \item $x_{soc}^* = 25$ y $\tau^* = 15$
\end{opts}

\q{En un mercado de autos usados, $50\%$ son de alta calidad (valorados en $20\,000$ por compradores y $16\,000$ por vendedores) y $50\%$ son de baja calidad (valorados en $8\,000$ por compradores y $6\,000$ por vendedores). Bajo información asimétrica, la disposición a pagar esperada inicial y el equilibrio final son:}
\begin{opts}
    \item $14\,000$; se comercian ambos tipos de autos al precio único de $14\,000$.
    \item $14\,000$; solo se comercian autos de baja calidad (los buenos salen).
    \item $16\,000$; se comercian ambos tipos de autos sin pérdida de bienestar.
    \item $12\,000$; el mercado colapsa totalmente y no se comercia ningún auto.
\end{opts}

\q{En un mercado con información asimétrica, los compradores valoran un bien bueno en $20\,000$ y uno malo en $8\,000$. Si el precio de reserva de los vendedores de bienes buenos es $16\,000$, la proporción mínima $q^*$ de bienes buenos necesaria para evitar su expulsión del mercado es:}
\begin{opts}
    \item $q^* = 1/2$
    \item $q^* = 3/5$
    \item $q^* = 2/3$
    \item $q^* = 3/4$
\end{opts}

\q{En un modelo de señalización laboral de Spence, la productividad de los trabajadores altos es $w_H = 100$ y la de los bajos es $w_L = 50$. El costo de adquirir $e$ años de educación es $c_H(e) = 10e$ para los altos y $c_L(e) = 25e$ para los bajos. El rango de niveles educativos $e^*$ que sostiene un equilibrio separador es:}
\begin{opts}
    \item $1 \le e^* \le 4$
    \item $4 \le e^* \le 5$
    \item $2 \le e^* \le 10$
    \item $2 \le e^* \le 5$
\end{opts}

\q{En un mercado laboral competitivo con información asimétrica, el $40\%$ de los postulantes son de alta productividad ($120$ soles/hora) y el $60\%$ son de baja productividad ($70$ soles/hora). En un equilibrio agrupador (pooling) donde nadie señaliza, el salario competitivo de equilibrio es:}
\begin{opts}
    \item $90$ soles/hora
    \item $85$ soles/hora
    \item $95$ soles/hora
    \item $100$ soles/hora
\end{opts}

\q{Una población asegurada enfrenta una pérdida potencial de $S/\ 10\,000$ y se divide en partes iguales entre individuos de alto riesgo ($p_H = 0.30$) y bajo riesgo ($p_L = 0.10$). Si una aseguradora competitiva no distingue tipos y ofrece seguro completo agrupador, la prima actuarialmente justa es:}
\begin{opts}
    \item $S/\ 1\,000$
    \item $S/\ 2\,000$
    \item $S/\ 1\,500$
    \item $S/\ 3\,000$
\end{opts}

\q{En un problema Principal-Agente con esfuerzo $e$ no observable, el ingreso esperado es $E[Y] = 20e$, el costo de esfuerzo del agente es $c(e) = 0.5e^2$ y su utilidad de reserva es $\bar{U} = 50$. Si el agente es neutral al riesgo, el esfuerzo óptimo $e^*$ y la renta fija $F^*$ que el principal le cobra por la franquicia son:}
\begin{opts}
    \item $e^* = 10$ y $F^* = 100$
    \item $e^* = 20$ y $F^* = 200$
    \item $e^* = 20$ y $F^* = 150$
    \item $e^* = 25$ y $F^* = 175$
\end{opts}

\q{Un agente neutral al riesgo elige entre esfuerzo alto (costo $20$, probabilidad $0.8$ de éxito) o esfuerzo bajo (costo $0$, probabilidad $0.3$ de éxito). Si el principal paga $w_H$ en caso de éxito y $w_L$ en caso de fracaso, el bono diferencial mínimo $(w_H - w_L)$ que induce esfuerzo alto es:}
\begin{opts}
    \item $w_H - w_L = 20$
    \item $w_H - w_L = 30$
    \item $w_H - w_L = 50$
    \item $w_H - w_L = 40$
\end{opts}

\q{Una empresa regulada puede tener costo marginal bajo ($c_L = 4$) o alto ($c_H = 8$), siendo su tipo información privada. Si el contrato regulatorio asigna al tipo de costo alto una producción $q_H = 25$ con transferencia que cubre exactamente su costo ($T_H = 200$), la renta informacional mínima que debe garantizarse al tipo eficiente para que diga la verdad es:}
\begin{opts}
    \item $R_L = 100$ soles
    \item $R_L = 50$ soles
    \item $R_L = 150$ soles
    \item $R_L = 200$ soles
\end{opts}

\q{Los vertimientos sin tratar de una ciudad ($200\,000\ m^3$/mes) generan un daño ambiental externo de $S/\ 1.50$ por $m^3$. Construir y operar una PTAR cuesta $S/\ 180\,000$ mensuales y elimina el $80\%$ del daño ambiental. El beneficio social neto mensual de ejecutar la PTAR es:}
\begin{opts}
    \item $S/\ 40\,000$ mensuales
    \item $S/\ 60\,000$ mensuales
    \item $S/\ 120\,000$ mensuales
    \item $S/\ 240\,000$ mensuales
\end{opts}
